% Compile from this directory with: pdflatex project_12.tex (twice).
\documentclass[11pt,a4paper]{article}
\usepackage[T1]{fontenc}
\usepackage[utf8]{inputenc}
\usepackage[margin=22mm]{geometry}
\usepackage{lmodern}
\usepackage{amsmath,amssymb}
\usepackage{enumitem}
\usepackage{xcolor}
\usepackage{microtype}
\usepackage{hyperref}
\definecolor{epflred}{HTML}{E3000B}
\hypersetup{colorlinks=true,urlcolor=epflred,linkcolor=epflred,
  pdftitle={ENG-654 Project 12: Nonsingular Posture Change on custom\_3R}}
\setlength{\parindent}{0pt}
\setlength{\parskip}{6pt}
\setlist[itemize]{leftmargin=*,itemsep=5pt,topsep=4pt}
\urlstyle{tt}

\begin{document}
{\small\textbf{EPFL \textbar\ ENG-654}\hfill Kinematics-Grounded Motion Planning for Robots}
\vspace{5mm}

{\color{epflred}\Large\bfseries Project 12}

{\LARGE\bfseries Nonsingular Posture Change on custom\_3R\par}

\textbf{Format:} Group of two students; simulation-based project.\\
\textbf{Course foundations:} Lectures 3, 5--7; Exercise 3.

\section*{Motivation}
For a cuspidal robot, two distinct configurations can reach the same
workspace position while belonging to one singularity-free configuration
region. A nonsingular connecting joint path then produces a closed tool path.
This explains why posture labels and endpoint positions alone are insufficient
to understand global motion planning.

\section*{Task}
\textbf{Overall objectives.} Use the exact custom\_3R model to construct and validate a candidate
nonsingular connection between two distinct IK solutions of one position,
and interpret both its joint-space path and its closed workspace image.

\begin{itemize}
  \item \textbf{Validate the exact model.} Read the URDF and establish its D--H
parameters, angle offsets, and selected operational point. Compare forward
kinematics and Jacobians with the course implementation. The exact robot is
orthogonal with non-intersecting consecutive axes; do not substitute the
idealized atlas dimensions.

  \item \textbf{Locate the IKs on reduced maps.} Plot $\det J_p$ and its zero set
in $(q_2,q_3)$, then map the critical curves and two-/four-IK regions in
$(\rho,z)$. Select a regular target and overlay all its distinct IKs. Identify
two candidate configurations in the same computed singularity-free component;
equal determinant signs alone are insufficient.

  \item \textbf{Construct the joint-space connection.} Use the Lecture 6
nonsingular-change demonstration to obtain a starting candidate, then construct
or adjust a piecewise-smooth path in the reduced component. Specify a continuous
$q_1$ interpolation connecting the full endpoint IKs. Recover and display the
complete three-joint trajectory.

  \item \textbf{Validate the workspace loop.} Compute the forward image of every
joint sample and check that the endpoint positions agree while endpoint
configurations are distinct modulo $2\pi$. Display both the three-dimensional
workspace loop and its meridian projection. Do not infer full Cartesian
closure from agreement in $(\rho,z)$ alone.

  \item \textbf{Test regularity and robustness.} Monitor $\det J_p$, the smallest
singular value, and optional joint-limit violations. Refine near the smallest
margin and compare the original path with one small path perturbation.
Separate strong sampled evidence from a continuous mathematical guarantee;
use an analytical bound only if readily available.
\end{itemize}

\newpage
\section*{Specifics and scope}
Study one target, one selected IK pair, and two candidate paths (the
original and one variation). The lecture's existing connection is an allowed
seed, but validation must use the exact model. The primary cuspidality study
uses periodic unconstrained joints; report separately whether the same path
respects the URDF limits. No obstacle model or general cusp-search algorithm
is required.

Use the position Jacobian $J_p$ for the 3R task.
Confirm that the meridian reduction is valid for the selected model and
operational point. Count distinct real IKs modulo periodicity, then filter
joint limits separately. Identify two-/four-IK regions where present and
explain any absent count rather than forcing both types into every map. Show both determinant values and the zero contour,
and distinguish sampled connectivity from a proof. Refine the decisive
region instead of claiming completeness from a coarse grid.

\section*{Expected outcome}
Explain cuspidality through a concrete nonsingular posture change, including
the relationship between reduced singularities, multiple IKs, and a closed
workspace image. Deliver the selected IK pair, reproducible joint paths,
paired plots, and a regularity/closure validation report.

Submit reproducible code, model parameters and test data, the required plots,
a concise 5--6-page report, and a short simulation demonstration. Explain
both successful cases and failures; conclusions must follow the numerical
and geometric evidence rather than an expected outcome.

\section*{Resources}
The exact custom\_3R is an orthogonal 3R robot with non-intersecting
consecutive joint axes. Its URDF is
\path{lectures_main/assets/models/custom_3R/custom_3R_new.urdf}; STL meshes are
in the same directory. Use the lecture website to inspect axes and frames,
verify D--H parameters, plot IKs, and animate configurations.
Use the exact-model nonsingular-change lab in Lecture 6 and the Lecture 3 IK
method. The idealized atlas can explain concepts but must not supply metric
results for this exact-URDF experiment.

Use the lecture website to verify D--H parameters, visualize IKs and
singularities, and simulate paths. All listed paths are relative to the
repository root. Start the website following \path{lectures_main/README.md}.

\section*{Workload and collaboration}
Budget approximately 60 hours per student (120 person-hours per pair)
for this 2-ECTS project, following EPFL's 30-hours-per-credit convention.\footnote{\href{https://www.epfl.ch/education/teaching/teaching-guide-2/getting-started/design-a-course_1/}{EPFL: Design a Course, workload guidance}.}
Deduct any lectures or exercises included in those credits. Allocate roughly
20 team-hours to model validation, 50 to the core work, 30 to experiments,
and 20 to reporting. Reuse the specified IK and simulations. Share validation
and interpretation even if modelling and implementation are divided between
students. Hardware execution is outside scope.

\end{document}
